\documentclass[11pt]{article}
\usepackage[margin=0.85in]{geometry}
\usepackage{amsmath,booktabs,hyperref,graphicx}
\hypersetup{colorlinks=true,urlcolor=blue}
\title{Turning-crossing collision: frame-level causal diagnosis}
\author{Pan Dynamics Collision Avoidance Research}
\date{October 2, 2026}
\begin{document}
\maketitle

\section{Finding and scope}
The 15 m/s turning-crossing collision was preceded by inaccurate affine
avoidance predictions, followed by persistently positive optimized safety
slack. It was not a solver timeout or an unobserved collision solely between
safe sample nodes. Both optimization stages returned exit flag 1 in the first
collision hold and the deepest-overlap hold. This is a diagnosis of controller
commit:
\begin{center}\small\texttt{2697350ca383bbbc326f2ca516129c740cc29fd0}\end{center}
No controller changes were made.

The final source-frozen campaign has exact state observations, a known target
trajectory, no road constraints, a 50 ms sample period, 16 primary intervals,
60 completion intervals, and a 0.10 m node buffer. Actual success requires
strictly positive body clearance, not a 0.10 m actual clearance. Recorded
states were fed back through the same controller to reconstruct the preceding
input plans. Two separate replays reproduced all first 35 issued inputs
exactly. An external diagnostic copy of the solver added only a save of the
selected primary point before the CLF solve; it reproduced the same commands.
All model comparisons below are offline and do not introduce an online gate.

\section{The collision frame}
Independent rectangle geometry applied to the original ODE45 dense replay
finds overlap from 1.638333 to 1.683333 s at the recorded sampling resolution.
The first collision occurs during the hold starting at 1.60 s (frame 33 when
the initial frame is numbered 1). The deepest sampled overlap occurs at
1.661667 s, during the 1.65 s hold (frame 34). Its signed separating-axis gap
is $-0.223768$ m. This is a geometric overlap measure, not crash deformation.

\begin{center}
\begin{tabular}{lrr}
\toprule
Quantity & First collision hold & Deepest-overlap hold\\
\midrule
Hold start [s] & 1.60 & 1.65\\
Signed gap at hold start [m] & 0.324356 & $-0.118889$\\
Primary objective & 0.21836258 & 0.21888935\\
Current-state slack lower bound & 0 & 0.21888931\\
First steering [rad] & 0.28581191 & 0.13262822\\
First signed longitudinal input & 0.22500019 & 0.47415941\\
Selected PCBF / CLF exit flags & 1 / 1 & 1 / 1\\
\bottomrule
\end{tabular}
\end{center}
The longitudinal input is named \texttt{brakingRatio}, but its sign convention
is positive drive and negative braking. The 1.65 s optimization starts with
overlap already present; its current-state violation cannot be undone by any
new control. Blaming that control alone would miss the earlier failure.
At 1.60 s the optimizer itself predicts overlap at 1.65 s: signed gap
$-0.118358$ m, versus $-0.118889$ m in the nonlinear rollout. Prediction
position error over this one hold is only 0.792 mm. The collision is therefore
already represented in the returned relaxed plan.

\section{How the plan lost safety}
For each selected frame, the complete optimized input sequence was rolled
out with the original nonlinear Fiala vehicle model from the same measured
state. The following compares vehicle rectangles at the same absolute future
time, 1.65 s. These are counterfactual open-loop plan evaluations, not the
closed-loop trajectory ultimately executed after further replanning.

\begin{center}
\begin{tabular}{rrrrr}
\toprule
Planning time & Affine gap & Nonlinear gap & Position error & PCBF optimum\\
{[s]} & {[m]} & {[m]} & {[m]} & \\
\midrule
0.75 & 0.127089 & 0.015114 & 0.089031 & $5.42\,10^{-7}$\\
0.80 & 0.112275 & $-0.069312$ & 0.132228 & $7.24\,10^{-9}$\\
0.85 & 0.231291 & $-0.250256$ & 0.889361 & $3.04\,10^{-7}$\\
0.90 & $-0.135776$ & $-0.205596$ & 0.244014 & 0.244982\\
1.60 & $-0.118358$ & $-0.118889$ & 0.000792 & 0.218363\\
\bottomrule
\end{tabular}
\end{center}
Gaps in this table are signed SAT gaps; negative entries mean actual rectangle
overlap. Thus by 0.80 s a zero-slack affine plan already does not represent a
safe nonlinear trajectory. At 0.85 s shifted and fresh-flow primary problems
both return infeasible; doubling the input correction box yields a near-zero
primary objective. The selected secondary plan has substantially worse
dynamic prediction consistency despite satisfying the affine safety model.

\subsection{The local model permits physically optimistic forces}
At the 0.85 s frame the selected affine trajectory has a peak normalized
combined axle-force magnitude of 1.594 in the exported first 24 stages:
\[
 \max_a\sqrt{(\widehat F_{y,a}/(\mu F_{z,a}))^2+b^2}=1.594.
\]
The nonlinear Fiala model limits the corresponding combined force to unity.
The local affine model has no explicit friction-circle rows and its tangent
can extrapolate outside this force envelope. The ratio already reaches
1.026 at 0.75 s and 1.063 at 0.80 s. These are predictions, not excessive
forces actually produced by the simulated plant.

The evidence identifies a concrete mechanism for optimistic dynamics; it
does not isolate tire-force extrapolation as the sole contributor to every
state error. Nonlinear propagation and orientation errors also matter.
Exact rectangle geometry at the affine state still shows positive clearance
at 0.80 and 0.85 s, whereas the nonlinear rollout overlaps. Consequently,
rectangle linearization error alone cannot explain those failures.
One-hold RK4 versus tight-tolerance ODE45 position discrepancies along the
selected 0.85 s plan reach only 0.052 mm over the evaluated prefix. This is
far smaller than the 0.889 m accumulated affine/nonlinear plan discrepancy.

\subsection{CLF priority protects the affine surrogate, not nonlinear safety}
The second optimization preserves the primary slack sum up to the configured
$10^{-6}$ tie tolerance, but can select quite different inputs among those
affine-equivalent plans. At 0.85 s the first input changes from
\[
 u_0^{\rm PCBF}=(0.251855,-0.493934),\qquad
 u_0^{\rm CLF}=(-0.021011,-0.100000).
\]
The latter both reverses the steering sign and weakens braking. The predicted
affine clearance at 1.65 s improves from 0.180568 to 0.231291 m, while the
affine/nonlinear position discrepancy worsens from 0.395739 to 0.889361 m.
Both raw primary and secondary nonlinear rollouts overlap at that future
node. At 0.80 s the primary has only 1.945 mm clearance at that particular
node and the secondary overlaps; this does not certify the primary's full
trajectory. The CLF stage is not skipped, and it does not violate the stated
affine lexicographic priority. The limitation is what that priority protects.

\subsection{A hard future constraint becomes soft before the conflict}
At planning time 0.85 s, the 1.65 s node is the start of stage 17, in the hard
completion tail. At planning time 0.90 s it is stage 16, in the slack-enabled
primary horizon. The controller can then return the same upcoming conflict
as a relaxed solution. At 0.90 s almost all the 0.244983 returned slack is
allocated to that stage; subsequent frames continue to have positive primary
objectives until the collision.

Fresh flow restarts occur, but they only change the local approximation and
its correction box. At 0.90--1.25 s the selected first steering corrections
reach approximately $+0.075$ rad and longitudinal corrections approximately
$-0.125$ relative to fresh-flow anchors. These are numerical trust bounds,
not physical steering constraints. In the measured 0.90 s case the shifted
problem is infeasible and the fresh one returns positive slack; the algorithm
does not expand a finite positive-slack solution or relinearize it to
convergence. It selects that result and proceeds to CLF optimization.
Merely minimizing each frame's relaxed sum does not enforce nonlinear
feasibility restoration, nor monotonic decrease of the violation across
changing models and the moving hard/soft boundary.

\section{Interpretation and limits}
The supported causal chain is: an overly optimistic local prediction supplies
near-zero modeled safety slack; the secondary stage can worsen the true
maneuver while keeping that modeled slack; later local problems openly admit
the upcoming collision through slack and fail to remove it before execution.
No road constraint, observation error, missed sampling node, or collision-frame
solver failure is needed to explain this run. Terminal infeasibility is not
established as the cause by this diagnostic.

The design issue is consistency between predicted safety and realizable
motion, followed by the absence of effective restoration once positive slack
appears. Stopping CLF optimization, adding a backup mode, or adding online
rejection checks is not demonstrated here as the necessary solution. A
revision should make the local dynamics credible over the allowed steps and
make restoration genuinely progress, while retaining the requested single
controller and PCBF-then-CLF ordering. This diagnosis does not prove global
unavoidability or global feasibility from any particular recorded state.

\section{Reproduction and validation}
The companion directory \path{report/COLLISION_FRAME_DIAGNOSIS_20261002}
contains scripts, exported nodes, selected original frames, primary/secondary
comparisons, hash manifest and column definitions. Both 35-frame replays
matched issued inputs exactly. Independent signed rectangle geometry
reproduced the previously reported collision interval and worst overlap.
The sole instrumentation patch writes diagnostic data before stage two.
No full scenario campaign or controller unit suite was rerun because the
production algorithm is unchanged. The initial diagnostic script's empty
struct append was corrected before the successful complete replay.
Original MATLAB snapshots and generated figures are retained outside the
repository and copied into the external research archive.
\end{document}
